\documentclass[11pt,a4paper]{article}
\usepackage[utf8]{inputenc}
\usepackage[T1]{fontenc}
\usepackage[spanish,es-nodecimaldot,es-tabla]{babel}
\usepackage[margin=2.5cm]{geometry}
\usepackage{lmodern}
\usepackage{microtype}
\usepackage{mathtools,amssymb}
\usepackage{booktabs,longtable,tabularx,array}
\usepackage{xurl}
\usepackage[hidelinks]{hyperref}
\hypersetup{pdftitle={Ecuación simplificada de producción y liberación de CO2},
  pdfsubject={Formulación, parámetros y fundamento bibliográfico}}
\urlstyle{same}
\newcolumntype{P}[1]{>{\raggedright\arraybackslash}p{#1}}
\newcommand{\qbio}{q_{\mathrm{bio}}}
\newcommand{\qprod}{q_{\mathrm{prod}}}
\newcommand{\qout}{q_{\mathrm{out}}}
\newcommand{\qpred}{q_{\mathrm{pred}}}
\newcommand{\qresp}{q_{\mathrm{resp}}}
\newcommand{\fferm}{f_{\mathrm{ferm}}}
\newcommand{\pos}[1]{\left[#1\right]_+}
\newcommand{\clip}{\operatorname{clip}}
\newcommand{\code}[1]{\texttt{\detokenize{#1}}}
\setlength{\parskip}{0.3em}
\setlength{\emergencystretch}{2em}
\renewcommand{\arraystretch}{1.15}
\setlength{\LTpre}{0.6em}
\setlength{\LTpost}{0.6em}
\allowdisplaybreaks[2]
\widowpenalty=10000
\clubpenalty=10000

\begin{document}
\begin{center}
  {\large\bfseries Ecuación simplificada de producción y liberación de CO\(_2\)}
\end{center}

\section{Alcance y notación}

Se documenta la variante evaluada en la comparación final de LAB019--LAB021:
el transitorio nutricional se incorpora dentro de la producción biológica
efectiva, se conservan la rampa y el factor de actividad nutricional, y se omite
el aporte aditivo de CO\(_2\) respiratorio. Se mantiene la modulación por oxígeno,
la activación inicial y la dinámica del depósito disuelto. La variante pertenece
al notebook de reformulación; no implica una sustitución de todos los modelos
de CO\(_2\) del repositorio.

La formulación combina un núcleo cinético con una capa fenomenológica de
producción y liberación. Sus fundamentos y las diferencias respecto de los
artículos revisados se explicitan en la sección~\ref{sec:bases}. No se atribuye
la ecuación completa a un único artículo.

\begin{longtable}{@{}P{2.0cm}P{12.5cm}@{}}
\caption{Notación de las tasas de CO\(_2\).}\label{tab:q}\\
\toprule Símbolo & Definición \\\midrule\endfirsthead
\toprule Símbolo & Definición \\\midrule\endhead
\bottomrule\endfoot
$q_{\mathrm{bio},0}$ & Producción biológica de la trayectoria contrafactual sin pulsos posteriores de N, antes de rampa, gain, transitorio, O\(_2\) y activación.\\
$q_{\mathrm{bio},N}$ & Producción biológica de la trayectoria cinética con los pulsos de N, antes de las correcciones de la capa de CO\(_2\).\\
$\qbio$ & Producción biológica efectiva con rampa, gain y transitorio, antes de O\(_2\) y activación. Esta es su definición en la variante simplificada.\\
$\qprod$ & Producción que ingresa al depósito disuelto, después de aplicar la activación y la modulación fermentativa por O\(_2\).\\
$\qout$ & Flujo físico liberado desde el depósito disuelto.\\
$\qpred$ & Flujo predicho que se compara con el sensor, después del factor de escala de matriz.\\
$q_{\mathrm{obs}}$ & Flujo observado válido de CO\(_2\).\\
\end{longtable}

Todas estas tasas se expresan en $\mathrm{g\,CO_2\,L^{-1}\,h^{-1}}$.
El tiempo $t$ se expresa en horas desde el origen temporal del experimento.
Se utiliza $[v]_+=\max(v,0)$ y
$\clip(v,0,1)=\min(\max(v,0),1)$. No es necesario introducir tasas adicionales
como $q_T$ o $q_{\mathrm{base}}$ para escribir esta variante.

\section{Ecuación principal}

La producción que alimenta el depósito disuelto es
\begin{equation}
\boxed{\qprod(t)=a(t)\,\fferm(t)\,\qbio(t).}
\label{eq:principal}
\end{equation}
La producción biológica efectiva se define por
\begin{equation}
\boxed{\begin{aligned}
\qbio(t)={}&[1+T_N(t)]\Bigl[\,[1+(g-1)r(t)]q_{\mathrm{bio},0}(t)\\
&\hspace{2.0cm}+r(t)\bigl(q_{\mathrm{bio},N}(t)-q_{\mathrm{bio},0}(t)\bigr)\,\Bigr]_+ .
\end{aligned}}
\label{eq:qbio}
\end{equation}
Aquí $a$ es la activación inicial, $\fferm$ la modulación fermentativa por
oxígeno, $r$ la rampa nutricional, $g$ el factor de actividad
\emph{gain}, y $T_N$ el incremento transitorio adimensional.
El transitorio multiplica la expresión completa de rampa y gain y se aplica
una sola vez dentro de $\qbio$.

\section{Definición química de las producciones biológicas}

\subsection{Trayectorias con y sin pulsos de nitrógeno}

Para $k\in\{0,N\}$ se utiliza
\begin{equation}
\boxed{q_{\mathrm{bio},k}(t)=Y_{\mathrm{CO_2/E}}X_k(t)
\bigl[\beta_{G0}\varphi_{G,k}(t)+\beta_{F0}\varphi_{F,k}(t)\bigr].}
\label{eq:raw}
\end{equation}
Las dos trayectorias tienen las mismas condiciones iniciales, temperatura y
theta. La trayectoria $0$ omite únicamente los pulsos posteriores de N; la
trayectoria $N$ conserva dichos pulsos. El subíndice $0$ no representa el
instante inicial ni implica una concentración inicial de N igual a cero.
Otros eventos del experimento se conservan en ambas simulaciones.

Los estados $X_k,G_k,F_k,E_k,N_k$ son, respectivamente, biomasa viable,
glucosa, fructosa, etanol y nitrógeno asimilable. Se expresan en
$\mathrm{g\,L^{-1}}$, equivalentes numéricamente a
$\mathrm{kg\,m^{-3}}$ en las variables cinéticas del código. Las masas
corresponden a la especie indicada por cada símbolo. La temperatura de entrada
$T(t)$ se expresa en grados Celsius y $T_K=T+273.15$ en kelvin.

\subsection{Glucosa, fructosa e inhibición por etanol}

Los factores cinéticos son
\begin{align}
\varphi_{G,k}(t)&=a_\beta(T)
\frac{G_k(t)}{G_k(t)+K_G(T)}
\frac{1}{1+i_E(T)E_k(t)},\label{eq:phiG}\\
\varphi_{F,k}(t)&=a_\beta(T)
\frac{F_k(t)}{F_k(t)+K_F(T)}
\frac{1}{1+i_G(T)G_k(t)}
\frac{1}{1+i_E(T)E_k(t)}.\label{eq:phiF}
\end{align}
Las fracciones con $G_k$ y $F_k$ representan la saturación por azúcar. El etanol
inhibe ambas contribuciones; la glucosa inhibe adicionalmente la contribución
de fructosa. La biomasa viable $X_k$ transforma las velocidades específicas en
tasas volumétricas.

Esta estructura corresponde a las expresiones separadas de glucosa y
fructosa de Zenteno et al.~\cite{zenteno2010}, ecuaciones (9)--(10), p.~41.
El artículo expresa las inhibiciones como $K_{IE}/(E+K_{IE})$ y
$K_{IG}/(G+K_{IG})$. Son equivalentes a las formas aquí utilizadas mediante
$i_E=1/K_{IE}$ e $i_G=1/K_{IG}$ a la misma temperatura. La relación de la
producción de etanol con $X$ y ambas contribuciones figura en su ecuación (12).

\subsection{Correcciones de temperatura}

Las funciones ejecutadas en el núcleo cinético son
\begin{align}
a_\beta(T)&=\exp\left[\frac{11000(T_K-296.15)}{296.15\,R\,T_K}\right],
\label{eq:abeta}\\
a_K(T)&=\exp\left[\frac{46055(T_K-293.15)}{293.15\,R\,T_K}\right],
\qquad R=8.314\ \mathrm{J\,mol^{-1}\,K^{-1}}.\label{eq:ak}
\end{align}
Los cinco factores térmicos de las constantes de saturación e inhibición
comparten actualmente la energía de activación y la temperatura de referencia,
por lo que se agrupan en $a_K$. Las constantes efectivas son
\begin{align}
K_G(T)&=\frac{\beta_{G0}}{s_G}a_K(T),&
K_F(T)&=\frac{\beta_{F0}}{s_F}a_K(T),\label{eq:sugarK}\\
i_G(T)&=\frac{i_G}{a_K(T)},&
i_E(T)&=\frac{i_E}{a_K(T)}.\label{eq:inhibition}
\end{align}
$i_G,i_E$ sin argumento son los parámetros de referencia de theta.
$K_G,K_F$ tienen unidades de concentración de su azúcar;
$i_G(T),i_E(T)$ tienen unidades de concentración inversa.
$\beta_{G0},\beta_{F0}$ se expresan en
$\mathrm{g\,E\,(g\,X)^{-1}\,h^{-1}}$. Las unidades de $s_G,s_F$ son las de
su $\beta$ divididas por concentración del azúcar respectivo, de modo que
los cocientes de la ecuación~\eqref{eq:sugarK} sean concentraciones.

Las dependencias de Arrhenius y las energías utilizadas están relacionadas con
Zenteno et al.~\cite{zenteno2010}, ecuaciones (5), (11), (13) y (21),
y tabla 2. La implementación usa $293.15$ K en $a_K$; el artículo escribe
$293.3$ K en su ecuación (21). Por ello se documenta una adaptación de la
estructura, no una reproducción literal de todos sus coeficientes.

\subsection{Coeficiente \texorpdfstring{CO\(_2\)}{CO2}/etanol y alcance estequiométrico}
\label{sec:yield}

El coeficiente conservado de la implementación es
\begin{equation}
Y_{\mathrm{CO_2/E}}=\frac{44.01}{2(46.07)}
=0.4776427176\ \mathrm{g\,CO_2\,(g\,E)^{-1}}.
\label{eq:yield}
\end{equation}
Es el factor que multiplica la producción de etanol en el código y se
mantiene para preservar la comparación numérica.

La reacción ideal de fermentación de una hexosa produce dos moles de etanol y
dos moles de CO\(_2\). Su cociente másico CO\(_2\)/etanol sería
$44.01/46.07\simeq0.955285$, el doble del valor de
\eqref{eq:yield}. El rendimiento CO\(_2\)/azúcar, del orden de
$2(44.01)/180.16\simeq0.489$, es otro cociente y no debe confundirse con
CO\(_2\)/etanol. Stroia y Lodin~\cite{stroia2025}, ecuación (8), p.~7,
emplean un rendimiento de CO\(_2\) por azúcar; David et al.~\cite{david2012},
ecuaciones (1)--(2), p.~354, relacionan aproximadamente las masas de etanol
y CO\(_2\), con una conversión empírica azúcar/CO\(_2\).
Ninguno de esos resultados justifica el factor $1/2$ adicional de
\eqref{eq:yield}. Se identifica, por tanto, como coeficiente operativo
heredado cuya interpretación estequiométrica requiere revisión separada.
No se modifica en esta simplificación ni se atribuye como tal a la bibliografía.

\section{Representación del efecto nutricional}

\subsection{Efecto cinético del nitrógeno añadido}

El N modifica la trayectoria $N$, especialmente el crecimiento y las historias
de consumo de azúcar y producción de etanol. La velocidad de crecimiento es
\begin{align}
\mu_k(t)&=\mu_0 a_\mu(T)\frac{N_k}{N_k+K_N(T)},\label{eq:mu}\\
a_\mu(T)&=\exp\left[\frac{59453(T_K-300)}{300\,R\,T_K}\right],&
K_N(T)&=\frac{\mu_0}{s_N}a_K(T).\label{eq:KN}
\end{align}
En cada evento nutricional de la trayectoria $N$ se aplica el salto exacto
\begin{equation}
N_N(t_j^+)=N_N(t_j^-)+\Delta N_j.
\label{eq:Npulse}
\end{equation}
La trayectoria $0$ omite esos saltos. La producción
$q_{\mathrm{bio},N}$ ya contiene el efecto cinético de la adición; no se
introduce otra saturación explícita por N en la ecuación~\eqref{eq:raw}.
El crecimiento limitado por N tiene respaldo en Zenteno et al.~\cite{zenteno2010},
ecuación (2), y Coleman et al.~\cite{coleman2007}, ecuación (6).

\subsection{Rampa de utilización y gain}

Se define
\begin{equation}
r(t)=\clip\left(\frac{t-t_p}{t_{\mathrm{rise}}},0,1\right),
\label{eq:ramp}
\end{equation}
donde $t_p$ es el primer evento nutricional y $t_{\mathrm{rise}}>0$ el tiempo
de rampa. Si no hay pulsos, $r=0$. La combinación antes del transitorio es
\begin{equation}
\pos{[1+(g-1)r]q_{\mathrm{bio},0}
+r(q_{\mathrm{bio},N}-q_{\mathrm{bio},0})}.
\label{eq:rampgain}
\end{equation}
Con $r=0$ se obtiene $q_{\mathrm{bio},0}$; con $r=1$ se obtiene
$[q_{\mathrm{bio},N}+(g-1)q_{\mathrm{bio},0}]_+$.
El gain modifica la actividad de la trayectoria contrafactual; no multiplica
por sí solo toda la producción de la trayectoria $N$. Con $g=0.25$ su
contribución basal se atenúa progresivamente. El operador de positividad es
necesario porque la combinación no constituye siempre una media convexa.

David et al.~\cite{david2012}, sección III-A, p.~356, distribuyen el N añadido
uniformemente durante $[t_{\mathrm{add}},t_{\mathrm{add}}+t_{\mathrm{rise}}]$.
Su utilización acumulada motiva una rampa lineal. Sin embargo, el artículo
aplica entradas distribuidas en balances de N y transportadores; aquí el N
cinético salta exactamente y la rampa combina dos trayectorias ya calculadas.
La expresión~\eqref{eq:rampgain} y su gain son una adaptación propia.
Además, el artículo hace depender el tiempo de respuesta y otros parámetros
del progreso fermentativo; el presente diagnóstico congela esos parámetros.

Con varios pulsos, la implementación mantiene una única rampa anclada al
primero. Todos los pulsos se incluyen en la trayectoria $N$ y en el transitorio.

\subsection{Transitorio nutricional}
\label{sec:transient}

Se conservan expresamente $A$ y la dosis de referencia $N_{\mathrm{ref}}$:
\begin{equation}
\boxed{T_N(t)=A\sum_j\frac{\Delta N_j}{N_{\mathrm{ref}}}K_j(t),
\qquad N_{\mathrm{ref}}=0.14\ \mathrm{g\,N\,L^{-1}}.}
\label{eq:TN}
\end{equation}
Para cada pulso,
\begin{equation}
K_j(t)=\frac{(1-e^{-u_j/\tau_r})e^{-u_j/\tau_d}}{\mathcal N_j},
\qquad u_j(t)=\max(t-t_j,0).
\label{eq:kernel}
\end{equation}
$A$ es adimensional; $\tau_r,\tau_d>0$ son tiempos en horas, y
$\mathcal N_j$ normaliza el máximo del kernel. El incremento comienza en cero,
sube y decae hasta cero. Para un pulso aislado con
$\Delta N_j=N_{\mathrm{ref}}$ cuyo kernel alcanza uno, el multiplicador máximo
es $1+A$.

La normalización antigua es el máximo del numerador de
\eqref{eq:kernel} en la grilla temporal del experimento. Depende, por tanto,
de la grilla y del horizonte. La nueva utiliza el máximo analítico:
\begin{equation}
u^*=\tau_r\ln\left(1+\frac{\tau_d}{\tau_r}\right),
\qquad \mathcal N_j=(1-e^{-u^*/\tau_r})e^{-u^*/\tau_d}.
\label{eq:peak}
\end{equation}
El tiempo de rampa $t_{\mathrm{rise}}$ y el tiempo de subida del transitorio
$\tau_r$ son parámetros diferentes.

\paragraph{Procedencia de la dosis de referencia.}
El valor $0.14\ \mathrm{g\,N\,L^{-1}}$ se tomó de la dosis considerada en el
diagnóstico histórico LAB016--LAB018. El protocolo asumido es 1.0 g de
SpringFerm Xtrem más 0.4 g de FDA en 2 L, usando la conversión del repositorio
de $0.20\ \mathrm{g\,N\,(g\,producto)^{-1}}$:
\begin{equation}
\Delta N_{\mathrm{protocolo}}
=\frac{(1.0+0.4)\,0.20}{2}
=0.14\ \mathrm{g\,N\,L^{-1}}.
\label{eq:doseprotocol}
\end{equation}
Es una dosis calculada bajo esos supuestos, no una medición directa de YAN,
una constante fisiológica, una dosis óptima ni una concentración de
semisaturación. No proviene de theta ni de los artículos revisados. Es una
normalización fija que permite interpretar $A$ por dosis de referencia.
Para otras fermentaciones se modifica $\Delta N_j$ y se conserva
$N_{\mathrm{ref}}$. La dosis debe corresponder a N asimilable por volumen,
no a la masa total del producto nutricional.

La formulación supone una amplitud transitoria proporcional a la dosis y
mantiene los tiempos característicos. Para $A=0.30$, dosis de 0.07, 0.08,
0.14 y 0.28 g N/L producen incrementos máximos aislados de 15\%, 17.14\%,
30\% y 60\%, respectivamente. Estos porcentajes describen el multiplicador
en $\qbio$, no directamente el pico del flujo del sensor. La extrapolación
lineal a otras dosis no ha sido validada con esta comparación.

\paragraph{Fundamento y límites.}
Malherbe et al.~\cite{malherbe2004} distinguen crecimiento y actividad
metabólica vinculada al N. David et al.~\cite{david2012}, sección III,
y Seguinot et al.~\cite{seguinot2018}, sección 3.1, muestran reactivación
fermentativa tras adiciones tardías, incluso con poco o ningún crecimiento
adicional. Esto respalda representar un efecto de actividad distinto del
incremento de biomasa. No se encontró en los nueve artículos revisados el
kernel exacto~\eqref{eq:kernel}, su normalización, el valor de
$N_{\mathrm{ref}}$ ni la ley universal de amplitud lineal~\eqref{eq:TN}.
Se documentan como hipótesis fenomenológicas propias.
Agrupar los mecanismos dentro de $\qbio$ simplifica la notación, pero conserva
su superposición y no prueba su identificabilidad independiente.

\section{Modulación fermentativa por oxígeno}

La función conservada es
\begin{equation}
\fferm(t)=f_{\min}+(1-f_{\min})
\frac{K_{\mathrm{ana}}^n}{K_{\mathrm{ana}}^n+O(t)^n},
\label{eq:fferm}
\end{equation}
con $f_{\min}=0.08$, $K_{\mathrm{ana}}=0.75\ \mathrm{mg\,O_2\,L^{-1}}$
y $n=2$. Para $O\to0$, $\fferm\to1$; con O\(_2\) elevado respecto de
$K_{\mathrm{ana}}$, $\fferm\to0.08$. Ese piso conserva actividad fermentativa
bajo condiciones oxigenadas; no es una tasa respiratoria.

El oxígeno disuelto $O$, en mg/L, satisface
\begin{align}
\frac{dO}{dt}&=-q_{O_2,\max}X_{\mathrm{eff}}
\frac{O}{K_O+O},& K_O&=0.25\ \mathrm{mg\,O_2\,L^{-1}},\label{eq:oxygen}\\
X_{\mathrm{eff}}&=\pos{X_0+r(X_N-X_0)},&
O(0)&=\alpha_O O_{\mathrm{sat}}(0).\label{eq:Xeff}
\end{align}
$q_{O_2,\max}$ tiene unidades de
$\mathrm{mg\,O_2\,(g\,X)^{-1}\,h^{-1}}$ y $\alpha_O$ es adimensional.
La correlación de inicialización utilizada es
\begin{equation}
O_{\mathrm{sat}}(t)=8.6e^{-0.024(T-20)}e^{-0.002E_N}
e^{-0.001(G_N+F_N)}\quad\mathrm{mg\,O_2\,L^{-1}}.
\label{eq:Osat}
\end{equation}
Los coeficientes de los exponentes se aplican a $T$ en grados Celsius y a
$E_N,G_N,F_N$ en g/L. No hay reposición de O\(_2\) en esta configuración.
Ni $g$ ni $1+T_N$ multiplican $X_{\mathrm{eff}}$ ni el consumo de oxígeno.

Stroia y Lodin~\cite{stroia2025}, ecuación (6), p.~6, presentan un modelo con
O\(_2\) y saturación del tipo Michaelis--Menten, pero con balances y
acoplamientos diferentes. Es un antecedente conceptual para incorporar O\(_2\),
no la fuente de la ley~\eqref{eq:fferm}, de sus valores 0.08 y 0.75, ni de
la correlación~\eqref{eq:Osat}. No se encontró respaldo directo para esas
expresiones exactas en el conjunto revisado.

\section{Activación inicial de la producción}

Se conserva una activación suave y acotada:
\begin{align}
a(t)&=3z(t)^2-2z(t)^3,\label{eq:gate}\\
z(t)&=\clip\left(\frac{t-t_{\mathrm{start}}}{D},0,1\right),\label{eq:z}\\
t_{\mathrm{start}}&=L+s(U-L),&
D&=\max\bigl(d(U-L),0.25\ \mathrm h\bigr).\label{eq:gateparams}
\end{align}
$s,d$ son las fracciones de comienzo y duración de la activación,
respectivamente. El intervalo $[L,U]$ se obtiene de dos modos:
\begin{enumerate}
\item En modo predicho, se utiliza el primer cruce de la trayectoria simulada
con pulsos en que $G_N+F_N$ disminuye 5 g/L o $E_N$ aumenta 2 g/L respecto
de sus valores iniciales. $U$ es el tiempo del cruce y $L$ el punto anterior
de la grilla. Si no hay cruce, el helper fija $L=U=t_0$.
\item En modo manual, se establece $[L,U]=[0,t_{\mathrm{onset}}]$.
Así, $t_{\mathrm{start}}=s\,t_{\mathrm{onset}}$ y
$D=\max(d\,t_{\mathrm{onset}},0.25\ \mathrm h)$.
\end{enumerate}
Si $U\le L+10^{-9}$ h, el código usa un escalón en $U$.
El onset manual condiciona un intervalo de activación; no desplaza directamente
toda la curva. Los valores efectivos dependen también de $s,d$.
La forma polinómica, los umbrales químicos y su parametrización son decisiones
del modelo del repositorio. No se identificó esta construcción exacta en
los artículos revisados. Las comparaciones con onset manual son diagnósticos
condicionados, no validaciones predictivas independientes del inicio.

\section{Depósito disuelto, liberación y predicción del sensor}

La concentración disuelta $C$, en g CO\(_2\)/L, se integra según
\begin{equation}
\frac{dC}{dt}=\qprod-\qout,\qquad C(0)=0.
\label{eq:pool}
\end{equation}
La ley de liberación continua es
\begin{equation}
\boxed{\qout=k_{\mathrm{release}}C
\left(\varepsilon_{\mathrm{release}}+
(1-\varepsilon_{\mathrm{release}})\frac{C}{C+C_{\mathrm{sat}}}\right),}
\label{eq:out}
\end{equation}
con $\varepsilon_{\mathrm{release}}=0.05$ y
$k_{\mathrm{release}}$ en h$^{-1}$. La solubilidad aparente es
\begin{equation}
C_{\mathrm{sat}}(t)=s_{\mathrm{CO_2}}\,
1.69e^{-0.032(T-20)}e^{0.0016E_N}e^{-0.0012(G_N+F_N)}
\quad\mathrm{g\,CO_2\,L^{-1}}.
\label{eq:Csat}
\end{equation}
$s_{\mathrm{CO_2}}$ es adimensional. La correlación se evalúa con $T$ en
grados Celsius y los estados químicos en g/L. En esta ley,
$C_{\mathrm{sat}}$ modula la fracción de liberación; no impone una cota rígida
a $C$. Finalmente,
\begin{equation}
\boxed{\qpred(t)=g_{\mathrm{matrix}}\qout(t).}
\label{eq:pred}
\end{equation}
El factor adimensional $g_{\mathrm{matrix}}$ escala la comparación con el
sensor. Es distinto de $g$, que interviene en la expresión nutricional
\eqref{eq:qbio}. El balance físico del depósito utiliza $\qout$, no
$\qpred$.

Zenteno et al.~\cite{zenteno2010}, sección 2.7, ecuaciones (14)--(16), p.~42,
distinguen CO\(_2\) producido, acumulación disuelta y flujo liberado. Ese es
el fundamento conceptual de separar producción y emisión. Su modelo utiliza
otra ley de transferencia, un cambio de régimen al saturarse y un aporte de
CO\(_2\) asociado al crecimiento que aquí no se incorpora explícitamente.
No se encontraron en los artículos revisados la ley racional continua
\eqref{eq:out}, su piso de 0.05, la correlación exacta~\eqref{eq:Csat}
ni el factor $g_{\mathrm{matrix}}$. Son cierres efectivos de la implementación,
y no una reproducción literal del modelo de Zenteno.

\section{Parámetros y procedencia}
\label{sec:params}

\subsection{Juegos de parámetros de la capa de \texorpdfstring{CO\(_2\)}{CO2}}

Las denominaciones \emph{antiguos} y \emph{nuevos} se refieren a la capa
de CO\(_2\), no a theta. Ambos juegos se pueden cruzar con ambos vectores
cinéticos. Los valores se presentan redondeados; el notebook conserva
la precisión original.

\begin{longtable}{@{}P{5.6cm}rrP{2.4cm}@{}}
\caption{Parámetros CO\(_2\) congelados.}\label{tab:co2params}\\
\toprule Parámetro & Antiguos & Nuevos & Unidad \\\midrule\endfirsthead
\toprule Parámetro & Antiguos & Nuevos & Unidad \\\midrule\endhead
\bottomrule\endfoot
$k_{\mathrm{release}}$ & 0.424588271353 & 0.096957891648 & h$^{-1}$\\
$s_{\mathrm{CO_2}}$ & 0.351389541107 & 2.50000000000 & Adimensional\\
$g_{\mathrm{matrix}}$ & 3.11575694438 & 1.44671770928 & Adimensional\\
$q_{O_2,\max}$ & 0.150000000000 & 0.253137377347 & mg O\(_2\)/(g X h)\\
$\alpha_O$ & 0.195399665651 & 0.179605552367 & Adimensional\\
$t_{\mathrm{rise}}$ & 64.6587037463 & 56.7418471260 & h\\
$g$ & 0.250000000000 & 0.250010532263 & Adimensional\\
$s$ & 0.735056695859 & 0.741574639341 & Adimensional\\
$d$ & 1.06390986814 & 0.457773307394 & Adimensional\\
$A$ & 0.30 & 1.50 & Adimensional\\
$\tau_r$ & 4.00 & 0.540335260310 & h\\
$\tau_d$ & 12.00 & 4.39925061955 & h\\
\end{longtable}

La normalización del kernel es máximo de grilla en el juego antiguo y máximo
analítico en el nuevo. Es una convención de cálculo, no un parámetro ajustado.
Las constantes restantes se mantienen: $N_{\mathrm{ref}}=0.14$ g N/L,
$K_O=0.25$ mg O\(_2\)/L, $K_{\mathrm{ana}}=0.75$ mg O\(_2\)/L,
$n=2$, $f_{\min}=0.08$ y $\varepsilon_{\mathrm{release}}=0.05$.

La correspondencia entre símbolos y nombres de configuración es
\begin{center}\small
\begin{tabular}{@{}ll@{}}
\toprule Símbolo & Nombre en el código\\\midrule
$k_{\mathrm{release}}$ & \code{kCO2_release_h}\\
$s_{\mathrm{CO_2}}$ & \code{CO2sat_scale}\\
$g_{\mathrm{matrix}}$ & \code{matrix_gain}\\
$q_{O_2,\max}$ & \code{O2_qmax_mg_gdw_h}\\
$\alpha_O$ & \code{O2_initial_scale}\\
$t_{\mathrm{rise}}$ & \code{pulse_t_rise_h}\\
$g$ & \code{pulse_activity_gain}\\
$s$ & \code{chem_activation_start_fraction}\\
$d$ & \code{chem_activation_duration_fraction}\\
$\tau_r$, $\tau_d$ & \code{tau_rise_h}, \code{tau_decay_h}\\
\bottomrule
\end{tabular}
\end{center}

\paragraph{Juego antiguo.}
Los nueve parámetros de O\(_2\), activación, rampa, solubilidad, liberación y
matriz provienen de la calibración natural histórica en
\path{laboratory_2026/results/co2_matrix_cross_validation_2026/fit_parameters.csv}.
Se añaden los valores manuales $A=0.30$, $\tau_r=4$ h y $\tau_d=12$ h
del diagnóstico histórico LAB016--LAB018. Esos tres valores no provienen
de theta ni de una calibración automática del transitorio.

\paragraph{Juego nuevo.}
Los seis parámetros de O\(_2\), activación y rampa se heredan del estado
\code{C_full_theta_refit}, matriz natural, en
\path{laboratory_2026/results/co2_matrix_cross_validation_2026_full_theta_sccm_corrected_no_lab010/three_state_co2_parameters.csv}.
Sobre ese estado se sustituyen $k_{\mathrm{release}},s_{\mathrm{CO_2}},
g_{\mathrm{matrix}},A,\tau_r,\tau_d$ por el snapshot de la curva azul
\emph{Actual con O2}, conservado como \code{BLUE_POOL_SNAPSHOT} en el notebook.
Esos seis valores proceden de un ajuste conjunto con theta\_full fijo,
LAB005/007/008 con peso 1 y LAB019/020/021 con peso 2. No corresponden al
ajuste diagnóstico marrón realizado sobre LAB016--LAB018. Leer solamente
el CSV del estado indicado no recupera el juego nuevo completo: se requieren
los seis reemplazos del snapshot.

\subsection{Theta antiguo y theta\_full}

Theta contiene 17 parámetros cinéticos. El theta antiguo efectivo conserva
once estimaciones naturales y completa seis con valores por defecto.
Theta\_full conserva esas once estimaciones y recupera los otros seis
del vector cinético conjunto. Los parámetros CO\(_2\) de la
tabla~\ref{tab:co2params}, incluido $A$, no forman parte de theta.

\begin{longtable}{@{}P{1.9cm}rrP{5.9cm}@{}}
\caption{Los seis parámetros cinéticos que cambian.}\label{tab:changedtheta}\\
\toprule Parámetro & Antiguo & Full & Función \\\midrule\endfirsthead
\toprule Parámetro & Antiguo & Full & Función \\\midrule\endhead
\bottomrule\endfoot
$s_N$ & 18 & 8.75144587094 & Determina $K_N=\mu_0a_K/s_N$.\\
$q_{XG}$ & 0.1125 & 0.081490986405 & Consumo de glucosa asociado al crecimiento.\\
$q_{XF}$ & 0.1125 & 0.070932706502 & Consumo de fructosa asociado al crecimiento.\\
$s_G$ & 0.03 & 0.097455098739 & Determina $K_G=\beta_{G0}a_K/s_G$.\\
$s_F$ & 0.03 & 0.178295205518 & Determina $K_F=\beta_{F0}a_K/s_F$.\\
$m_0$ & 0.01 & 0.018338579393 & Consumo específico de azúcar por mantenimiento.\\
\end{longtable}

\begin{longtable}{@{}P{1.9cm}rP{8.5cm}@{}}
\caption{Los once parámetros cinéticos comunes.}\label{tab:commontheta}\\
\toprule Parámetro & Valor común & Función \\\midrule\endfirsthead
\toprule Parámetro & Valor común & Función \\\midrule\endhead
\bottomrule\endfoot
$\mu_0$ & 0.0787279167845 & Escala de crecimiento.\\
$q_N$ & 0.0210803824285 & Consumo de N asociado al crecimiento; no es una tasa de CO\(_2\).\\
$\beta_{G0}$ & 0.428878002997 & Producción específica de etanol asociada a glucosa.\\
$\beta_{F0}$ & 3.44039917477 & Producción específica de etanol asociada a fructosa.\\
$q_{EG}$ & 1.71280679377 & Consumo fermentativo específico de glucosa.\\
$q_{EF}$ & 5.02301675017 & Consumo fermentativo específico de fructosa.\\
$i_G$ & 0.056037437270 & Inhibición de fructosa por glucosa.\\
$i_E$ & 0.040167940980 & Inhibición por etanol.\\
$K_{d0}$ & 0.000150249548 & Escala de muerte celular.\\
$\gamma_{G0}$ & 0.115558791368 & Producción específica de glicerol asociada a glucosa.\\
$\gamma_{F0}$ & 0.000100001622 & Producción específica de glicerol asociada a fructosa.\\
\end{longtable}

$\mu_0,K_{d0}$ tienen unidades h$^{-1}$; $s_N$ tiene unidades
$\mathrm{L\,(g\,N)^{-1}\,h^{-1}}$. Los parámetros $q_N$, $q_{XG}$, $q_{XF}$,
$q_{EG}$, $q_{EF}$ y $m_0$ son velocidades específicas de consumo de la especie
correspondiente por masa de biomasa y hora. Los $\gamma$ son velocidades
específicas de producción de glicerol por masa de biomasa y hora.
Las unidades de $s_G,s_F,i_G,i_E$ se indican en las ecuaciones
\eqref{eq:sugarK}--\eqref{eq:inhibition}.

Las fuentes son
\path{laboratory_2026/results/estimability_historical_natural/theta.csv}
para las once estimaciones naturales,
\path{laboratory_2026/results/estimability_historical_natural/theta_natural_full.csv}
para el vector full consolidado, y
\path{shared/results/new_must_glycerol_overnight_validation/theta_multistart_07.csv}
para los seis valores recuperados. Los seis valores antiguos proceden de
\code{DEFAULT_THETA} del módulo cinético.

Con iguales $\beta_{G0},\beta_{F0}$, el aumento de $s_G,s_F$ reduce
$K_G,K_F$, aumentando las fracciones de saturación a igual concentración de
azúcar. Esto puede elevar la producción y acelerar el consumo, contribuyendo
a picos mayores y colas más cortas. Los cambios en crecimiento y mantenimiento
también modifican las trayectorias; el efecto final no depende de un único
parámetro.

\section{Generación de las trayectorias cinéticas}
\label{sec:kinetics}

Para definir completamente los estados utilizados en $q_{\mathrm{bio},k}$,
se presenta el sistema cinético que permanece congelado. Con
$\psi_k=a_\mu N_k/(N_k+K_N)$ y las $\varphi$ de
\eqref{eq:phiG}--\eqref{eq:phiF}, entre eventos se integra
\begin{align}
\dot X_k&=(\mu_0\psi_k-k_{d,k})X_k,&
\dot X_{d,k}&=k_{d,k}X_k,\label{eq:biomass}\\
\dot N_k&=-q_N\psi_kX_k,\label{eq:N}\\
\dot G_k&=-X_k\left[q_{XG}\psi_k+q_{EG}\varphi_{G,k}
+m(T)\chi_k\frac{G_k}{G_k+F_k}\right],\label{eq:G}\\
\dot F_k&=-X_k\left[q_{XF}\psi_k+q_{EF}\varphi_{F,k}
+m(T)\chi_k\frac{F_k}{G_k+F_k}\right],\label{eq:F}\\
\dot E_k&=X_k[\beta_{G0}\varphi_{G,k}+\beta_{F0}\varphi_{F,k}],\label{eq:E}\\
\dot{Gly}_k&=X_k[\gamma_{G0}\varphi_{G,k}+\gamma_{F0}\varphi_{F,k}].\label{eq:Gly}
\end{align}
$X_d$ es biomasa muerta y $Gly$ glicerol, ambos en g/L. No aparecen
directamente en la producción instantánea de CO\(_2\). El mantenimiento es
\begin{equation}
m(T)=m_0\exp\left[\frac{37681(T_K-293.3)}{293.3\,R\,T_K}\right],
\qquad \chi_k=\frac{G_k+F_k}{G_k+F_k+S_m},\quad S_m=1\ \mathrm{g\,L^{-1}}.
\label{eq:maint}
\end{equation}
El consumo de azúcares incluye crecimiento, fermentación y mantenimiento.
La estructura tiene una relación directa con Zenteno et al.~\cite{zenteno2010},
ecuaciones (6)--(13); la parametrización mediante velocidades $q$ y pendientes
$s$, el factor de disponibilidad $\chi_k$ y las ecuaciones de glicerol son
adaptaciones de la implementación. La dependencia del crecimiento y la
actividad respecto de N y temperatura también está respaldada conceptualmente
por Coleman et al.~\cite{coleman2007}, sin reproducir sus regresiones térmicas.

La muerte celular conserva la ecuación de Zenteno et al.~\cite{zenteno2010},
ecuaciones (3)--(4), p.~41:
\begin{align}
T_d(E_k)&=-0.0001E_k^3+0.0049E_k^2-0.1279E_k+315.89,\label{eq:Td}\\
k_{d,k}&=\begin{cases}
0,&T_K<T_d(E_k),\\[1mm]
K_{d0}\exp\!\left[C_{de}E_k+
\dfrac{E_{td}(T_K-305.65)}{305.65\,R\,T_K}\right],&T_K\ge T_d(E_k).
\end{cases}\label{eq:death}
\end{align}
$C_{de}=0.0415\ \mathrm{L\,(g\,E)^{-1}}$,
$E_{td}=130000\ \mathrm{J\,mol^{-1}}$ y $T_d$ se expresa en kelvin.
El término $C_{de}E_k$ está fuera de la fracción térmica, como en el código
y en el artículo.

El núcleo cinético genera primero los estados. La capa CO\(_2\) usa esas
trayectorias sin retroalimentarlas: activación, modulación por O\(_2\),
rampa, gain y transitorio no modifican las ecuaciones cinéticas de esta
sección. En particular, $\qprod$ no representa aquí una reconstrucción
estequiométrica completa de todos los balances metabólicos.

\section{Implementación numérica y alcance de la simplificación}

\subsection{Convenciones de integración}

Las expresiones analíticas omiten regularizaciones pequeñas para facilitar
la lectura. La implementación evalúa estados no negativos, añade
$\epsilon=10^{-8}$ a los denominadores cinéticos pertinentes y utiliza
$i_G/(a_K+\epsilon)$, $i_E/(a_K+\epsilon)$. En mantenimiento emplea
$G+F+\epsilon$ tanto para las proporciones como para la disponibilidad.
Los pulsos se integran como saltos exactos en sus tiempos registrados.

En la capa CO\(_2\) se utiliza una grilla de 0.25 h. El O\(_2\) se actualiza
con Euler y se restringe a valores no negativos. Para el depósito se aplica
un límite de CO\(_2\) disponible:
\begin{align}
q_{\mathrm{out},i}=\min\Biggl\{&k_{\mathrm{release}}C_{i-1}
\left(0.05+0.95\frac{C_{i-1}}
{\max(C_{i-1}+C_{\mathrm{sat},i},10^{-12})}\right),\nonumber\\
&\frac{C_{i-1}}{\Delta t_i}+q_{\mathrm{prod},i-1}\Biggr\},\label{eq:Eulerq}\\
C_i&=\max\bigl[C_{i-1}+\Delta t_i(q_{\mathrm{prod},i-1}-q_{\mathrm{out},i}),0\bigr].\label{eq:EulerC}
\end{align}
Se inicializan $C_0=0$ y $q_{\mathrm{out},0}=0$. Estas convenciones se conservan para
las tres variantes comparadas; el depósito se reintegra para cada producción.
Las condiciones iniciales, temperatura, tiempos y dosis de nutrición y onset
manual son entradas separadas de theta y de los juegos CO\(_2\).

\subsection{Diferencia respecto de la referencia original}

Denótese por $\mathcal R$ únicamente la expresión común de rampa y gain
de~\eqref{eq:rampgain}; es una abreviatura algebraica, no una tasa adicional.
La referencia original y la variante simplificada son
\begin{align}
\qprod^{\mathrm{original}}&=a(1+T_N)[\fferm\mathcal R+\qresp],\label{eq:original}\\
\qbio^{\mathrm{simplificada}}&=(1+T_N)\mathcal R,&
\qprod^{\mathrm{simplificada}}&=a\fferm\qbio^{\mathrm{simplificada}}.\label{eq:simplified}
\end{align}
En la referencia,
\begin{equation}
\qresp=\frac{44.01}{32000}\,
q_{O_2,\max}X_{\mathrm{eff}}\frac{O}{K_O+O}.
\label{eq:oldresp}
\end{equation}
Mover el transitorio a $\qbio$ conservando respiración retira solamente
$aT_N\qresp$. Omitir además respiración retira de la referencia
$a(1+T_N)\qresp$. La parte fermentativa se conserva algebraicamente.
Retirar $\qresp$ no elimina $O$ ni $\fferm$.

En LAB019--LAB021, con todos los parámetros CO\(_2\) antiguos, onset manual
28/32/24 h y ambos theta, se verificaron por separado las producciones,
integraciones del depósito y curvas graficadas. Mover el transitorio produjo
diferencias máximas en $\qpred$ de hasta
$6.7\times10^{-16}\ \mathrm{g\,L^{-1}\,h^{-1}}$; omitir también respiración
produjo una diferencia máxima de
$3.52\times10^{-6}\ \mathrm{g\,L^{-1}\,h^{-1}}$.
La diferencia máxima de cada curva relativa a su pico fue inferior a
0.00042\%. Esto explica la superposición visual en esa configuración.
El resultado respalda la simplificación en los casos evaluados, no una
equivalencia universal con respiración distinta de cero.

\section{Fundamento bibliográfico y decisiones propias}
\label{sec:bases}

Se revisaron los nueve PDF disponibles en
\path{context/modeling}, dentro de \path{fermentation_model}.
Las referencias se asignan a componentes concretos; el respaldo conceptual
no se considera validación de todos los coeficientes ni de la ecuación completa.

\begin{longtable}{@{}P{3.2cm}P{11.3cm}@{}}
\caption{Correspondencia entre componentes y bibliografía revisada.}\\
\toprule Referencia & Relación con la formulación \\\midrule\endfirsthead
\toprule Referencia & Relación con la formulación \\\midrule\endhead
\bottomrule\endfoot
Zenteno et al. (2010)\newline\cite{zenteno2010} &
Respaldo directo de saturación separada de G/F, inhibición por E y G,
crecimiento limitado por N, mantenimiento y muerte térmica.
Respaldo conceptual del depósito disuelto y la distinción producción/emisión.
Su ley de liberación y su producción de CO\(_2\) difieren de las actuales.\\
Malherbe et al. (2004)\newline\cite{malherbe2004} &
Respaldo de la distinción entre crecimiento y actividad metabólica asociada
al N; transporte de azúcar e influencia de temperatura y etanol.
No contiene nuestra combinación de trayectorias ni el kernel exacto.\\
David et al. (2012)\newline\cite{david2012} &
Respaldo del tiempo finito de utilización del N y de la reactivación mediante
transportadores, particularmente con adición tardía.
Nuestra rampa se inspira en esa utilización distribuida, pero su ubicación
algebraica y el transitorio son diferentes.\\
Seguinot et al. (2018)\newline\cite{seguinot2018} &
Evidencia experimental de aumento de actividad tras N en fase estacionaria
sin incremento de población, sección 3.1, figura 1 y tabla 2.
No identifica nuestra ley de dosis ni los tiempos del kernel.\\
Coleman et al. (2007)\newline\cite{coleman2007} &
Respaldo del crecimiento limitado por N y de la dependencia de actividad,
temperatura e inactivación. Usa otras variables y dependencias térmicas;
no se le atribuyen los coeficientes Arrhenius presentes.\\
Stroia y Lodin (2025)\newline\cite{stroia2025} &
Antecedentes para usar la tasa CO\(_2\) como variable de seguimiento/control,
para modelos con O\(_2\) y para conversiones azúcar/CO\(_2\).
No respalda exactamente nuestro factor Hill, depósito, kernel ni conversión
operativa CO\(_2\)/etanol.\\
Colombié et al. (2007)\newline\cite{colombie2007} &
Relaciona tasa CO\(_2\) y potencia de fermentación, sección 2.2,
ecuación (6), y trata propiedades térmicas del mosto.
Contexto físico; no es fuente de la ecuación nutricional ni de las correlaciones
de solubilidad presentes. El PDF se llama \code{Colombie_2006.pdf}, pero
la publicación en volumen es de 2007.\\
Martínez et al. (1999)\newline\cite{martinez1999} &
Antecedente de seguimiento y control de fermentación mediante CO\(_2\) y
temperatura, secciones 1--2 y 4.
No proporciona los bloques específicos de la variante simplificada.\\
Fernández-Fernández y Pérez-Correa (2007)\newline\cite{fernandez2007} &
Trata producción de CO\(_2\), consumo de O\(_2\) y mantenimiento en un
biorreactor de fermentación en sustrato sólido, sección 2.2.1,
ecuaciones (4)--(5). Es un sistema diferente; no se usa como fundamento
directo para omitir respiración ni para parametrizar fermentación de mosto.\\
\end{longtable}

Se consideran decisiones propias de esta implementación: la combinación
contrafactual de rampa y gain; el kernel normalizado de subida y caída; la
referencia de dosis y su escalamiento lineal; el factor fermentativo Hill y
sus constantes; la activación polinómica y sus umbrales; las correlaciones
efectivas de O\(_2\) y CO\(_2\); la liberación racional continua y la escala de
matriz. Su elección puede ser compatible con los fenómenos descritos en la
literatura, pero las expresiones y valores exactos no se verificaron en esos
artículos. La omisión de CO\(_2\) respiratorio se fundamenta en la comparación
numérica indicada, no en una afirmación bibliográfica de ausencia de respiración.

\appendix
\section{Trazabilidad de la implementación}

Las rutas siguientes son relativas a \path{fermentation_model}. El artefacto
de comparación es
\path{laboratory_2026/notebooks/co2_reformulation/02_nutrition_structure_comparison.source.ipynb},
con su versión ejecutada
\path{laboratory_2026/notebooks/co2_reformulation/02_nutrition_structure_comparison.executed.ipynb}.
La figura final evalúa las tres ubicaciones/omisiones del transitorio y la
respiración mediante el helper
\path{laboratory_2026/notebooks/co2_reformulation/_nutrition_structure_comparison.py},
funciones \code{compare_transient_location} y \code{pool_from_production}.

El núcleo cinético es
\path{shared/run_new_must_glycerol_estimability_doe.py}, funciones
\code{kinetic_terms}, \code{rhs} y \code{simulate}.
La capa original de CO\(_2\) es
\path{laboratory_2026/run_co2_matrix_cross_validation_2026.py}; las funciones
\code{effective_qprod_grid} y \code{_bounded_smoothstep_activation} definen
O\(_2\), rampa/gain y activación. El cálculo de la producción biológica y las
correlaciones efectivas se encuentran en
\path{pilot_2025/run_pilot_2025_co2_solubility_integrated_doe.py}.
El factor de conversión~\eqref{eq:yield} se define como
\code{CO2_G_PER_G_ETHANOL} en
\path{shared/run_secondary_joint_campaign_doe.py}.
El diagnóstico histórico que fija el protocolo de la dosis de referencia es
\path{laboratory_2026/notebooks/lab016_018_natural_must_holdout.ipynb}.
La redacción de este documento no cambia esos módulos, notebooks ni parámetros.

\begin{thebibliography}{9}
\small
\setlength{\itemsep}{0.2em}
\setlength{\parskip}{0pt}
\bibitem{zenteno2010}
Zenteno, M.~I., Pérez-Correa, J.~R., Gelmi, C.~A., y Agosin, E. (2010).
Modeling temperature gradients in wine fermentation tanks.
\emph{Journal of Food Engineering}, 99, 40--48.
\url{https://doi.org/10.1016/j.jfoodeng.2010.01.033}.

\bibitem{malherbe2004}
Malherbe, S., Fromion, V., Hilgert, N., y Sablayrolles, J.-M. (2004).
Modeling the effects of assimilable nitrogen and temperature on fermentation
kinetics in enological conditions.
\emph{Biotechnology and Bioengineering}, 86(3), 261--272.
\url{https://doi.org/10.1002/bit.20075}.

\bibitem{david2012}
David, R., Dochain, D., Mouret, J.-R., Vande Wouwer, A., y Sablayrolles, J.-M.
(2012). Modeling of the boost effect originated by nitrogen addition during
wine-making fermentation. En \emph{2012 20th Mediterranean Conference on
Control \& Automation (MED)}, Barcelona, 3--6 de julio, pp.~353--358. IEEE.

\bibitem{seguinot2018}
Seguinot, P., Rollero, S., Sanchez, I., Sablayrolles, J.-M., Ortiz-Julien, A.,
Camarasa, C., y Mouret, J.-R. (2018). Impact of the timing and the nature of
nitrogen additions on the production kinetics of fermentative aromas by
\emph{Saccharomyces cerevisiae} during winemaking fermentation in synthetic
media. \emph{Food Microbiology}, 76, 29--39.
\url{https://doi.org/10.1016/j.fm.2018.04.005}.

\bibitem{coleman2007}
Coleman, M.~C., Fish, R., y Block, D.~E. (2007).
Temperature-dependent kinetic model for nitrogen-limited wine fermentations.
\emph{Applied and Environmental Microbiology}, 73(18), 5875--5884.
\url{https://doi.org/10.1128/AEM.00670-07}.

\bibitem{stroia2025}
Stroia, N., y Lodin, A. (2025). Modeling of must fermentation processes for
enabling CO\(_2\) rate-based control. \emph{Mathematics}, 13, 1653.
\url{https://doi.org/10.3390/math13101653}.

\bibitem{colombie2007}
Colombié, S., Malherbe, S., y Sablayrolles, J.-M. (2007).
Modeling of heat transfer in tanks during wine-making fermentation.
\emph{Food Control}, 18, 953--960.
\url{https://doi.org/10.1016/j.foodcont.2006.05.016}.

\bibitem{martinez1999}
Martínez, G., López, A., Esnoz, A., Vírseda, P., e Ibarrola, J. (1999).
A new fuzzy control system for white wine fermentation.
\emph{Food Control}, 10, 175--180.

\bibitem{fernandez2007}
Fernández-Fernández, M., y Pérez-Correa, J.~R. (2007).
Realistic model of a solid substrate fermentation packed-bed pilot bioreactor.
\emph{Process Biochemistry}, 42, 224--234.
\url{https://doi.org/10.1016/j.procbio.2006.08.003}.
\end{thebibliography}

\end{document}
